\clearpage
\VerdierPhysicalPageAnchor{0229}%
\begingroup\small\setlength{\parskip}{0pt}\setlength{\emergencystretch}{1.5em}
\setlength{\abovedisplayskip}{4pt plus 2pt minus 2pt}%
\setlength{\belowdisplayskip}{4pt plus 2pt minus 2pt}%

\VerdierSemanticAnchor{ch4d-4-1-6-cont}
\noindent complexe $n$-uple défini par :

\VerdierSemanticAnchor{ch4d-4-1-6-1}
\begin{equation}
({}^{m}\!X(p,q))^\sigma=
\begin{cases}
X^\sigma,&-q\leq\sigma_m<-p,\\
0,&\sigma_m<-q\ \lor\ -p\leq\sigma_m,
\end{cases},
\tag{4.1.6.1}
\end{equation}

\noindent les différentielles de ${}^{m}\!X(p,q)$ sont les différentielles de $X$
sauf peut-être si la source ou le but de la différentielle considérée sont
des objets nuls.

Pour tout couple $(p,q)$ et $(p',q')$ d'éléments de
$\mathcal{F}\ell_1(\widetilde{\mathbb Z})$, tels que $p\leq p'$ et
$q\leq q'$, désignons par :

\VerdierSemanticAnchor{ch4d-4-1-6-2}
\begin{equation}
f\bigl((p,q),(p',q')\bigr)\colon
{}^{m}\!X(p,q)\longrightarrow{}^{m}\!X(p',q')
\tag{4.1.6.2}
\end{equation}

le morphisme de complexes $n$-uples dont les composants sont des
morphismes identiques sauf peut-être si la source ou le but sont des objets
nuls. Enfin, pour tout élément $(p,q,r)$ de
$\mathcal{F}\ell_2(\widetilde{\mathbb Z})$, désignons par :

\VerdierSemanticAnchor{ch4d-4-1-6-3}
\begin{equation}
g(p,q,r)\colon{}^{m}\!X(q,r)
\longrightarrow T_m\bigl({}^{m}\!X(p,q)\bigr)
\tag{4.1.6.3}
\end{equation}

($T_m$ désigne le foncteur de translation dans la direction $m$
(chap.~I, 2.3.1) le morphisme dont les composants sont :

\VerdierSemanticAnchor{ch4d-4-1-6-4}
\begin{equation}
\left\{
\begin{aligned}
g(p,q,r)^\sigma&=0,
  &&\sigma_m\neq-q-1,\\
g(p,q,r)^\sigma&=-d_X^{\sigma,m},
  &&\sigma_m=-q-1.
\end{aligned}
\right.
\tag{4.1.6.4}
\end{equation}

Introduisons alors de nouvelles notations. Tout d'abord, désignons, par
abus de notation, par
$\int_n\colon\operatorname{comp}^n(\mathcal A)\longrightarrow
\operatorname{comp}(\mathcal A)$ la $\Pi$-extension ou la
$\Sigma$-extension du foncteur complexe simple associé (sous réserve
d'existence de produits ou sommes dénombrables dans $\mathcal A$
(chap.~I, 2.2.3). Posons alors :

\VerdierSemanticAnchor{ch4d-4-1-6-5}
\begin{equation}
\left\{
\begin{aligned}
\int_n{}^{m}\!X(p,q)
  &={}^{m}\!Y(p,q),\\
\int_n{}^{m}\!X(-\infty,q)
  &={}^{m}\!Y(-\infty,q)={}^{m}\!Y(q),\\
\int_n{}^{m}\!X(-\infty,+\infty)=\int_n X
  &={}^{m}\!Y(-\infty,+\infty)={}^{m}\!Y(+\infty)=Y.
\end{aligned}
\right.
\tag{4.1.6.5}
\end{equation}
\endgroup
